% Clase 6 --- Ángulo sólido: la imagen del globo (§2.4).
% Se dibuja una mancha sobre un globo esférico y se lo infla: el ÁREA crece
% como r^2 pero Omega = S/r^2 no cambia. La esfera entera da 4 pi.
\begin{center}
\begin{tikzpicture}[scale=1]

  % globo chico
  \begin{scope}
    \draw[figline] (0,0) circle (1.15);
    \draw[figaux] (0,0) ellipse (1.15 and 0.38);
    \fill[figred, opacity=0.75] (52:1.15) arc (52:8:1.15)
      .. controls (0.85,0.35) and (0.7,0.72) .. (52:1.15);
    \node[figlblaux, right=3pt] at (1.2,0.55) {$S$};
    \draw[figvecaux] (0,0) -- (215:1.15);
    \node[figlblaux, anchor=north east] at (215:0.85) {$r$};
    \node[figlbl, align=center] at (0,-1.85) {globo};
  \end{scope}

  % flecha "se infla"
  \draw[figvecres] (2.0,0) -- (3.1,0);
  \node[figlbl, figred, above=2pt] at (2.55,0.02) {se infla};

  % globo grande: la misma mancha, proyectada
  \begin{scope}[xshift=5.5cm]
    \draw[figline] (0,0) circle (1.85);
    \draw[figaux] (0,0) ellipse (1.85 and 0.61);
    \fill[figred, opacity=0.75] (52:1.85) arc (52:8:1.85)
      .. controls (1.37,0.56) and (1.13,1.16) .. (52:1.85);
    \node[figlblaux, right=3pt] at (1.95,0.9) {$S'$};
    \draw[figvecaux] (0,0) -- (215:1.85);
    \node[figlblaux, anchor=north east] at (215:1.35) {$r'$};
    \node[figlbl, align=center] at (0,-2.55) {mismo cono, mismo $\Omega$};
  \end{scope}

  % resultado
  \node[figlbl, align=center] at (2.75,-3.35)
    {$\Omega = \dfrac{S}{r^{2}} = \dfrac{S'}{r'^{2}}$};

\end{tikzpicture}\\[0.5em]
{\small\itshape Al inflar el globo la mancha se agranda ---el \textbf{área}
cambia--- pero el \textbf{ángulo sólido} no, porque todas las áreas crecen como
$r^{2}$. Es exactamente la cancelación que hace cero el flujo del cascarón
cónico, y la razón por la que en la ley de Gauss aparece un $4\pi$.}
\end{center}
